%=====================================================================
\chapter{Philosophy of Science for Computer Scientists}\label{ch:philosophy}
%=====================================================================
\locator{1}{Part~1 --- Foundations of Inquiry in Computing}

\begin{outcomes}
By the end of this chapter you will be able to:
\begin{tight}
  \item \textbf{Define} ontology, epistemology and methodology, and explain how
    a choice at one level constrains the next.
  \item \textbf{Locate} a computing study within five research paradigms from
    its methods alone.
  \item \textbf{Distinguish} deduction, induction and abduction, and identify
    which one design science actually uses.
  \item \textbf{Justify} your own paradigm choice in the terms a Departmental
    Research Committee will use.
  \item \textbf{Recognise} the two paradigms already encoded in the IUB
    dissertation quality criteria.
\end{tight}
\end{outcomes}

\begin{prereq}
\chref{contribution} --- particularly the idea that different contribution
types owe different evidence. This chapter explains \emph{why} they differ.
\end{prereq}

\begin{keyterms}
ontology $\bullet$ epistemology $\bullet$ axiology $\bullet$ methodology
$\bullet$ paradigm $\bullet$ positivism $\bullet$ post-positivism $\bullet$
interpretivism $\bullet$ critical realism $\bullet$ pragmatism $\bullet$
deduction $\bullet$ induction $\bullet$ abduction $\bullet$ falsification
$\bullet$ demarcation
\end{keyterms}

\section{Why This Chapter Is Not Optional}

Computing scholars tend to regard philosophy of science as decoration --- a
paragraph to insert in Chapter Three of the thesis, copied from a textbook,
declaring the study ``positivist'' because that word appeared in a paper the
supervisor liked. Examiners read a great many such paragraphs and are not
fooled by any of them.

There is a practical reason to take this seriously, and it has nothing to do
with sounding scholarly. \textbf{Your methods already commit you to a
philosophical position, whether or not you state one.} If you run a controlled
experiment, you have assumed that the effect you are measuring exists
independently of your measuring it, and that it is stable enough to reappear
for the next investigator. If you conduct unstructured interviews and build
categories from them, you have assumed something quite different --- that the
phenomenon of interest is partly constituted by how participants understand it.
Those assumptions are not compatible, and a thesis that runs both without
noticing produces a Chapter Three that contradicts Chapter Five.

\begin{definitionbox}[The four questions, in order]
\begin{tight}
  \item \term{Ontology} --- what kinds of thing exist? Is ``technical debt'' a
    real property of a codebase, or a metaphor practitioners use to talk about
    their discomfort?
  \item \term{Epistemology} --- what can count as knowledge of those things,
    and what is the relationship between the knower and the known? Can I know
    about technical debt without participating in the team that has it?
  \item \term{Axiology} --- what is the role of values? Should the researcher
    be neutral, or is neutrality itself a position?
  \item \term{Methodology} --- given the answers above, what procedures are
    legitimate?
\end{tight}
\medskip
The order matters: each answer constrains the next. Choosing a method first and
back-filling a philosophy is the standard error, and it is visible.
\end{definitionbox}

\section{Five Paradigms, and What They Look Like in Computing}

\begin{center}
\small
\begin{longtable}{@{}L{2.5cm}L{3.5cm}L{3.5cm}L{4.2cm}@{}}
\toprule
\textbf{Paradigm} & \textbf{What exists} & \textbf{What counts as knowledge} &
\textbf{Typical computing work} \\
\midrule
\endfirsthead
\toprule
\textbf{Paradigm} & \textbf{What exists} & \textbf{What counts as knowledge} &
\textbf{Typical computing work} \\
\midrule
\endhead
\bottomrule
\endfoot
Positivism
  & A single reality, independent of observers, governed by regularities
  & Verified observation; law-like generalisation
  & Algorithm complexity; performance benchmarking; anything with a
    ground truth that nobody disputes \\
Post-positivism
  & A single reality, but knowable only imperfectly and probabilistically
  & Conjectures that survive attempts at refutation; results always provisional
  & Most empirical software engineering and HCI experimentation; the working
    position of \chref{experiments} and \chref{inference} \\
Interpretivism / constructivism
  & Multiple realities, constructed by participants
  & Understanding of meaning in context; the researcher is part of the
    instrument
  & Studies of how teams adopt tools; why a process fails in one organisation
    and works in another (\chref{qualitative}) \\
Critical realism
  & A real world with underlying causal mechanisms, only partly observable
  & Explanation by mechanism: what generative structure would produce what we
    see?
  & Case study research seeking \emph{why} an intervention worked here and not
    there (\chref{casestudy}) \\
Pragmatism
  & Whatever is useful for the problem; the ontological question is deferred
  & What works, judged against a purpose
  & Design science (\chref{designscience}); most mixed-methods work
    (\chref{mixed}) \\
\end{longtable}
\end{center}

\begin{conceptbox}[The honest position of most computing research]
Very little computing research is strictly positivist, though a great many
theses claim to be. Positivism proper holds that we can verify universal laws
by observation --- a position largely abandoned in the philosophy of science by
the mid-twentieth century.

\medskip
What most quantitative computing research actually practises is
\textbf{post-positivism}: there is a real effect; our measurement of it is
imperfect; we advance by proposing something that could be refuted and failing
to refute it; every conclusion carries an interval rather than a full stop.
That is the position underlying null hypothesis testing, confidence intervals,
replication, and the entire apparatus of \chref{validity}.

\medskip
If you write ``this study adopts a positivist paradigm'' and then report a
confidence interval, you have described one philosophy and practised another.
Write post-positivist, and mean it.
\end{conceptbox}

\section{Three Ways to Reason, and the One Nobody Teaches}

\begin{center}
\small
\begin{tabular}{@{}L{2.3cm}L{5.0cm}L{6.6cm}@{}}
\toprule
\textbf{Mode} & \textbf{Moves from} & \textbf{In computing research} \\
\midrule
Deduction
  & General rule $\rightarrow$ particular prediction
  & Theory says caching helps read-heavy workloads; therefore this workload
    should speed up; test it. Conclusion is certain \emph{if} premises hold \\
Induction
  & Particular observations $\rightarrow$ general rule
  & Forty repositories show the pattern; perhaps it holds generally. Conclusion
    is probable, never certain, and the next repository can refute it \\
Abduction
  & Surprising observation $\rightarrow$ best explanation of it
  & Latency spikes every 40 minutes; what mechanism would produce exactly that?
    Generates the hypothesis that the other two then test \\
\bottomrule
\end{tabular}
\end{center}

Abduction is the mode that actually starts most research, and it is almost
never named in methods chapters. You notice something that does not fit, and
you reach for an explanation that would make it unsurprising. That explanation
is a conjecture, and it earns nothing until deduction turns it into a
prediction and evidence tests the prediction.

\begin{conceptbox}[Design science is abductive, not inductive]
This matters for how you write up design science work (\chref{designscience}).
A designer confronting a problem does not induce a solution from observations;
they conjecture an artefact whose properties, \emph{if} the conjecture about
the problem is right, would resolve it. The artefact is a hypothesis in
executable form. Evaluation is the deductive test of it.

\medskip
Writing this honestly --- ``we conjectured that the bottleneck was X, designed
Y on that assumption, and evaluated whether Y removed the bottleneck'' --- is
far more convincing than the usual retrospective narrative in which the correct
design appears fully formed and the evaluation confirms it.
\end{conceptbox}

\section{Falsification, and Its Limits}

Popper's answer to the problem of induction is the intellectual backbone of
\chref{contribution}'s falsifier line: no number of confirming observations
establishes a universal claim, but one genuine counter-example refutes it, so
science advances by proposing bold conjectures and trying hard to break
them~\cite{popper2002logic}.

This is the right working attitude, and it is why ``what would falsify this?''
is the most useful question in the course. But three qualifications matter for
a computing researcher.

\textbf{Falsification is rarely clean in practice.} When an experiment
contradicts a prediction, you may reject the hypothesis --- or the measurement
instrument, the sample, the implementation, or an auxiliary assumption. This is
the Duhem--Quine problem, and it is entirely familiar to anyone who has had a
benchmark disagree with a paper and spent a week finding the configuration
difference. Rejecting a claim honestly requires ruling out those alternatives,
which is precisely what a threats-to-validity section does.

\textbf{Fields do not abandon a framework on one refutation.} Kuhn's account of
normal science, anomaly and paradigm shift~\cite{kuhn1962structure} describes
what actually happens: anomalies accumulate and are tolerated until an
alternative framework is available. You will see this in your own literature.
A result that contradicts a dominant approach is not, on its own, going to
displace it, and a dissertation that claims otherwise will be met with
scepticism.

\textbf{Not all good computing research is falsificationist at all.}
An impossibility proof is not refuted by evidence; a qualitative study
generating a grounded theory is not trying to test a conjecture; an artefact
demonstrating feasibility establishes an existence claim. Falsification is the
right discipline for the empirical traditions, not a universal test of
scholarly worth.

\section{Is Computing a Science?}

The demarcation question --- what separates science from non-science --- lands
awkwardly on our discipline, and it is worth having a position because
examiners occasionally ask.

Computing is at least three things at once. Where it proves theorems about
computability and complexity, it behaves like mathematics: knowledge is a
priori and evidence is proof. Where it builds systems to specification, it
behaves like engineering: success is fitness for purpose. Where it studies what
happens when people, organisations and machines interact --- which is most of
software engineering, information systems and HCI --- it behaves like an
empirical science, with all the messiness that implies.

The practical consequence is that \emph{you must say which one your study is},
because the standard of evidence differs. Confusion here is a real source of
examiner disagreement: a formal researcher may find an empirical study
unrigorous because it proves nothing, while an empiricist may find a formal
result irrelevant because it was never confronted with data. Both are applying
their own tradition's standard, and both are wrong to apply it universally.
The \emph{Four Traditions} box in every chapter of this book exists to keep
that distinction visible.

\begin{regbox}[Two paradigms, written into the regulations --- PhD \S14.3]
The regulations do not use the word paradigm, but they encode the distinction
exactly, by giving \emph{separate} quality criteria for the two kinds of work.

\medskip
\textbf{\S14.3(A), qualitative research} asks whether the research
``illuminate[s] the subjective meaning, actions and contexts of those being
researched''; whether the design was adapted responsively to real-life
settings; whether the description is detailed enough for a reader to interpret
meaning and context; and whether ``subjective perceptions and experiences
[are] treated as knowledge in their own right''.

\medskip
\textbf{\S14.3(B), quantitative research} asks about reliability
(``are the results repeatable?''), validity (``does it measure what it
says?''), internal validity, external validity --- both ecological and
population --- and replicability.

\medskip
\S14.3(C) then requires \emph{appropriateness of the methods to the aims of the
study}. That clause is the one that penalises a mismatch between your stated
paradigm and your actual practice.
\end{regbox}

\begin{pitfall}[Three ways the philosophy chapter goes wrong]
\textbf{Paradigm as decoration.} A paragraph asserting a paradigm, followed by
a study that ignores it. The test: delete the paragraph. If nothing about the
rest of the thesis would have to change, it was decoration.

\medskip
\textbf{Method-first, philosophy-backfilled.} Choosing a survey because surveys
are familiar, then justifying it philosophically afterwards. Detectable because
the justification never rules anything out --- it is compatible with every
method the student might have chosen.

\medskip
\textbf{``Mixed methods'' as a way of avoiding the question.} Combining
qualitative and quantitative work is legitimate and often excellent
(\chref{mixed}). But pragmatism is a considered position about deferring
ontological commitment in favour of what serves the problem --- not a licence
to hold incompatible assumptions simultaneously without noticing. A good mixed
methods study says explicitly what each strand is for and how the two are
brought together.
\end{pitfall}

\begin{worked}[Locating a study from its methods alone]
A paper reports: semi-structured interviews with 14 developers at three firms,
transcripts coded openly then axially, categories member-checked with
participants, no statistical testing, findings presented as five themes with
extended quotations.

\medskip
\textbf{Ontology.} Multiple constructed realities --- the study is interested
in how these developers understand their situation, not in a rate or an effect.
\textbf{Epistemology.} Knowledge is co-produced; the interviewer is part of the
instrument, which is why member checking matters.
\textbf{Axiology.} Values are acknowledged rather than eliminated; the
researcher's position is disclosed.
\textbf{Paradigm.} Interpretivist.

\medskip
Now notice what would be a category error: asking whether these 14 developers
are ``representative'' and computing a margin of error. Statistical
representativeness is a positivist criterion imported into a study that never
claimed it. The right questions are whether the sample produced the kind of
knowledge needed to understand the structures at work, whether the description
is thick enough to interpret, and whether saturation was reached
(\chref{qualitative}) --- which is, almost word for word,
\reg{PhD Regs 2024}{\S14.3(A)}.
\end{worked}

\begin{traditions}{philosophical commitment}
\tradEMP{Post-positivist by default. A real effect exists; our estimate of it
is uncertain; we report the uncertainty and expect to be replicated.}
\tradDSR{Pragmatist. Truth is deferred in favour of utility for a purpose;
the artefact is an abductive conjecture and evaluation is its test.}
\tradQUAL{Interpretivist or critical realist. Either meaning is constructed and
must be understood in context, or real mechanisms underlie observed events and
the task is to explain them.}
\tradFORM{Rationalist. Knowledge is a priori, evidence is proof, and empirical
observation is irrelevant to the truth of the theorem --- though very relevant
to whether anyone should care about it.}
\end{traditions}

\begin{lab}[Read the philosophy off three papers]
Choose three papers in your area that use visibly different methods --- ideally
one experiment, one qualitative study, one formal or design paper.

\medskip
For each, \emph{without reading any explicit statement of paradigm}, infer from
the methods alone: what the authors take to exist, what they accept as
evidence, and what role they give themselves. Then find whatever paradigm claim
the paper makes, and compare.

\medskip
You will frequently find either no statement at all, or one that does not match
the practice. Note which, and why you think it happened. Bring the mismatch to
the seminar.
\end{lab}

\begin{checkpoint}
\begin{tightnum}
  \item A study measures the effect of code review on defect density across 200
    projects and reports a 95\% confidence interval. Which paradigm is being
    practised, and why is it not, strictly, positivism?
  \item Explain the Duhem--Quine problem using an example from benchmarking.
  \item Which mode of reasoning generates a hypothesis, and which two test it?
    Why is the generating one so rarely described in methods chapters?
  \item Your supervisor says the study must be ``objective''. Under
    interpretivism, what does that request mean, and what would you offer
    instead?
\end{tightnum}
\end{checkpoint}

\begin{chaptersummary}
\begin{tight}
  \item Your methods commit you to a philosophy whether you state one or not.
    The purpose of stating it is to keep the thesis internally consistent.
  \item Ontology constrains epistemology, which constrains methodology. Choose
    in that order; back-filling is visible to examiners.
  \item Five paradigms cover nearly all computing research. Most quantitative
    computing work is post-positivist, not positivist, and should say so.
  \item Abduction generates hypotheses; deduction and induction test them.
    Design science is abductive, and reads better when written that way.
  \item Falsification is the right discipline for empirical work but is
    complicated by the Duhem--Quine problem and does not apply to formal or
    purely exploratory work.
  \item \reg{PhD Regs 2024}{\S14.3} already encodes the qualitative /
    quantitative split as two separate quality checklists, and \S14.3(C)
    penalises a mismatch between the paradigm you claim and the one you use.
\end{tight}
\end{chaptersummary}

\begin{reviewq}
  \item A thesis states a positivist paradigm, uses a survey instrument
    validated in another country without re-validation, and reports
    significance tests without effect sizes. Identify the philosophical
    inconsistency and the two methodological consequences that follow from it.
  \item Critical realism is sometimes proposed as the natural paradigm for
    software engineering case study research. Construct the argument for that
    claim, then give the strongest objection to it.
  \item Is a formal proof of an algorithm's correctness a scientific
    contribution under Popper's demarcation criterion? Does your answer
    embarrass the criterion, or the question?
  \item \reg{PhD Regs 2024}{\S14.3(A)} treats ``subjective perceptions and
    experiences as knowledge in their own right''. Reconcile this with
    \S14.3(B)'s demand for replicability, given that a single dissertation may
    contain both strands.
\end{reviewq}
